Mostly fixing errors in 4qClassification.lean

The intended refactor (pasting from chat with Claude):

"Yes, that's exactly right, and you should feel confident about it. Let me state it precisely to make sure we're aligned:

  You have φ : ↥K →* MulAut ↥Q (abstract Sylow subgroups). You want the theorem output to mention CyclicGroup q and CyclicGroup 4. So before calling classify_sdp, you use eQ : ↥Q ≃* CyclicGroup q and
  h_K_C4.some : ↥K ≃* CyclicGroup 4 to conjugate φ into φ' : CyclicGroup 4 →* MulAut (CyclicGroup q), and get the bridge iso ↥Q ⋊ φ ≃* CyclicGroup q ⋊ φ'.

  Now classify_sdp is fed φ' operating on the concrete types (CyclicGroup q, CyclicGroup 4). The output says CyclicGroup q ⋊ φ' ≃* CyclicGroup q ⋊ (sdpCanonicalAction r for CyclicGroup q, CyclicGroup 4).

  On the "exactly equal vs extremely close" question: if canonicalC4OnCqAction is defined as sdpCanonicalAction applied to (CyclicGroup q, CyclicGroup 4) with exactly the same hN and hK that classify_sdp
  uses, they are definitionally equal by congruence with no unfolding — same function, same explicit arguments. So it's "exactly equal" from Lean's perspective, and the r=1 proof step is literally just
  chaining the three isos together.

  The key insight you had (that I missed): the transport step belongs before classify_sdp, not after. That eliminates the need for sdpCanonicalAction_transport entirely."

"Yes, that's exactly what I mean. canonicalC4OnCqAction should literally be defined as sdpCanonicalAction instantiated at the concrete cyclic groups:

  noncomputable def canonicalC4OnCqAction {q : ℕ} [hq : Fact q.Prime] (h : q ≠ 2) :
      CyclicGroup 4 →* MulAut (CyclicGroup q) :=
    haveI : Fact (Nat.Prime 2) := ⟨by norm_num⟩
    sdpCanonicalAction (show 2 ≠ q from by omega) h 2 1 Nat.one_pos
      (by rw [card_cyclicGroup, pow_one])   -- Nat.card (CyclicGroup q) = q^1
      (by rw [card_cyclicGroup]; norm_num)  -- Nat.card (CyclicGroup 4) = 2^2
      1 (one_le_min_two_factorization_two h)

  The "exactly the same" part doesn't mean the proof terms need to be literally identical text. By proof irrelevance, any two proofs of Nat.card (CyclicGroup q) = q^1 are definitionally equal, so it doesn't
  matter if one comes from by rw [card_cyclicGroup, pow_one] and the other from by simp [card_cyclicGroup]. They're the same.

  The reason this is fast and the earlier attempt wasn't: when Lean sees sdpCanonicalAction ... (CyclicGroup q) (CyclicGroup 4) ... and canonicalC4OnCqAction ... (which unfolds to the exact same
  sdpCanonicalAction ... (CyclicGroup q) (CyclicGroup 4) ...), it only needs to check argument equality by congruence — same head function, same N, same K, all proof args equal by proof irrelevance. It never
  unfolds sdpCanonicalAction's body (which is where Classical.choice and the expensive stuff lives).

  The earlier timeout happened because I was comparing sdpCanonicalAction applied to (↥Q, ↥K) with sdpCanonicalAction applied to (CyclicGroup q, CyclicGroup 4) — different N and K type arguments — so Lean
  couldn't use congruence and was forced to unfold the body to check equality."

Claude gave up on me mid refactor so I want you to finish the job.

Essentially, the thing is, if you check git history, you would see use of horrible horrible bridges because the semi direct product between a cyclic p and q group classification result was stated in a bad way, which gave rose to the need for a lot of adaptors/bridges. Now its been improved and so it should be quite nice. Could you please make the code for 4q compile and make sure its neat and clean?

Here's the math (rough notes):

Lemma 1: For groups of order $p^2 q$, there must exist either a normal Sylow $p$-group or a normal Sylow $q$-group.

Lemma 2: If $K$ is a cyclic $p$-group, two homomorphisms $f, g : K → \text{Aut}(H)$ define
isomorphic semi-direct products $H \rtimes K$ if they have the same image.

Lemma 3: Semi-direct products $C_{q^n} ⋊ C_{p^m}$ ($q$ odd prime, $p ≠ q$) are classified up to
isomorphism by $r ∈ \{0, …, \min(m, d)\}$ where $d = v_p(q - 1)$, giving $\min(m, d) + 1$ classes. Every action $f$ belongs to exactly one class.

Lemma 4: If $|K| = p$, then two homomorphisms $f, g : K → \text{Aut}(H)$ define isomorphic semi-direct products $H \rtimes K$ if their images are conjugate subgroups of $\text{Aut}(H)$.

# Case 1: There exists a normal Sylow $p$-group

$|N| = p^2$ and by Schur-zassenhaus theorem there exists a complementary subgroup $U$ of $G$ with $|U| = q$ such that $G \cong N \rtimes U$. We have $U \cong C_q$.

## Case 1.1: $N \cong C_{p^2}$
$\text{Aut}(N) = C_{p(p-1)}$
Let $\phi : C_q \rightarrow C_{p(p-1)}$. Done by Lemma 3 when $p$ is odd. When $p = 2$, we get $\phi : C_q \rightarrow C_2$ so only trivial homomorphism is possible as $q \nmid 2$ for odd $q$.

(For $p = 2$, only trivial homomorphism)
(For $q = 2$, then we get 2 homomorphisms)

$C_{p^2 q}$ always
$C_{p^2} \rtimes C_q$ if $q \mid p - 1$.
## Case 1.2: $N \cong C_p \times C_p$
$\text{Aut}(N) \cong \text{GL}_2(p)$
Let $\phi : C_q \rightarrow \text{GL}_2(p)$

???
### Special case $p = 2$
$GL_2(2) \cong S_3$
Either $\phi$ is a trivial action, or its image is a subgroup of order $q$ in $S_3$, hence $q = 3$. Only one subgroup of order $3$ so gives one group.

### Special case $q = 2$
$\text{Im}(\phi)$ is trivial or $\langle A \rangle$ where $A \in GL_2(p)$ is of order $2$. By lemma 4 we only need to look at similarity classes and pick a representative from each class. Then $A^2 = 1$ implies minimal polynomial of $A$ divides $x^2 - 1$ and since $p$ is odd, it has distinct factors hence diagonalisable. Hence each such matrix similar to a diagonal matrix $\text{diag}(\alpha, \beta)$ such that $\alpha^2 = 1, \beta^2 = 1$ hence $\alpha = \pm 1, \beta = \pm 1$ ($t^2 - 1$ has at most $2$ roots because we are over a field). So $\alpha = \beta = 1$ is order $1$ so we don't consider it. We can check there are two similarity classes $\alpha = 1, \beta = -1$ (or switch them) and one for $\alpha = \beta = -1$.

Side note: Lemma 4 not needed because the general condition $\phi_2(u) = \alpha \circ \phi_1(\beta(u)) \circ \alpha^{-1}$ becomes equivalent to similarity because $\text{Aut}(C_q) = C_1$ as $q = 2$ so $\beta = \text{id}$ and that $\phi_1, \phi_2$ are determined by where they send the generator as $U$ is cyclic.

C_3 -> Aut(G)

# Case 2: There exists a normal Sylow $q$-group

$|N| = q$ and by Schur-zassenhaus theorem there exists a complementary subgroup $U$ of $G$ with $|U| = p^2$ such that $G \cong N \rtimes U$. We have $N \cong C_q$. $\text{Aut}(N) \cong C_{q-1}$.

## Case 2.1 $U \cong C_{p^2}$
Let $\phi : C_{p^2} \rightarrow C_{q-1}$. Done by Lemma 3 when $q$ is odd. Only trivial homomorphism possible when $q = 2$.

(When $p = 2$, have to see if $q$ is $1$ mod $4$)
(When $q = 2$, we get only trivial homomorphism)

$C_{p^2q}$ always
$C_q \rtimes_1 C_{p^2}$ if $v_p(q - 1) \geq 1$
$C_q \rtimes_2 C_{p^2}$ if $v_p(q - 1) \geq 2$
## Case 2.2 $U \cong C_p \times C_p$
Let $\phi : C_p \times C_p \rightarrow C_{q-1}$

$\text{Im}(\phi) \cong (C_p \times C_p) / \ker(\phi)$ so order of image is either $1, p$ or $p^2$. If it is $1$, then we have the trivial homomorphism giving $G \cong C_p \times C_p \times C_{q}$. If it is $p^2$, then we must have $\ker(\phi) = 1$ but $\text{Im}(\phi)$ would then be a non-cyclic subgroup of the cyclic group $C_{q-1}$, impossible. If it is $p$, then ?????

(When $p = 2$, one trivial and one non-trivial?)
(When $q = 2$, we get only trivial homomorphism)

$C_p \times C_p \times C_q$
$C_q \rtimes (C_p \times C_p)$ if $p \mid q - 1$